\documentclass[11pt]{amsart}
\usepackage[margin=1.15in]{geometry}
\usepackage{amssymb,hyperref,tikz}
\newtheorem{theorem}{Theorem}
\newtheorem{lemma}[theorem]{Lemma}
\theoremstyle{remark}\newtheorem*{remark}{Remark}
\newcommand{\st}{\,:\,}
\title{The superdiagonals of the Mahonian triangle}
\author{Vamshi Jandhyala}
\date{25 November 2026}
\begin{document}
\begin{abstract}
Claesson, Frankl\'in and Steingr\'imsson showed that the generating functions of the subdiagonals of the Mahonian
triangle are $S_i(x)=(xC(x))^iS_0(x)$, and conjectured that the superdiagonals satisfy
$x^iS_{-i}(x)=S_0(x)C(x)^{-i}+P_i(x)/(1-x)^i$ with $P_i\in\mathbb Z[x]$ of degree $3i-2$. We prove this, with an
explicit formula for $P_i$. The argument is Lagrange inversion for Laurent series together with the observation that
the product defining the Mahonian numbers differs from Euler's product by a factor
$\prod_{j>n}(1-q^j)^{-1}=1+q^{n+1}/(1-q)+O(q^{2n+2})$.
\end{abstract}
\maketitle

\section{The statement}
How many permutations of $[n]$ have exactly $n$ inversions, or $n+1$, or $n-1$? Along any fixed diagonal $k=n-i$ of
the Mahonian triangle (Figure~\ref{fig:mech}) these numbers form a single sequence, and Claesson, Frankl\'in and
Steingr\'imsson \cite{CFS} found that all the diagonals below the main one are governed by one series: each is the
main diagonal times a power of $xC(x)$. They conjectured that the diagonals above it are almost as simple. This note
proves that conjecture.

Let $I_n(k)$ be the number of permutations of $[n]$ with $k$ inversions, so that
$\sum_k I_n(k)q^k=\prod_{j=1}^{n}(1-q^j)/(1-q)$, and for every integer $i$ put
\[
S_i(x)=\sum_{n\ge0}I_n(n-i)\,x^n ,
\]
with $I_0(0)=1$ and $I_n(k)=0$ when no permutation has $k$ inversions. All identities below are identities of formal
series, and empty sums are zero.
Let $C(x)=(1-\sqrt{1-4x})/(2x)$ be the Catalan generating function. In \cite[Theorem~3]{CFS} it is shown that
$S_i=(xC)^iS_0$ for $i\ge0$, and \cite[Conjecture~15]{CFS} (repeated in Frankl\'in's thesis \cite{F}) asserts that for
every $i\ge1$
\[
S_{-i}(x)\,x^i=S_0(x)\,C(x)^{-i}+\frac{P_i(x)}{(1-x)^i}
\qquad\text{for a polynomial }P_i\in\mathbb Z[x]\text{ of degree }3i-2 .
\tag{1}
\]
In words: above the main diagonal, the series $x^iS_{-i}$ is the subdiagonal formula continued to negative $i$, up to
a rational correction whose only pole is at $x=1$. For example $P_1=-1+2x$ and $P_2=-1+4x-4x^2+x^3+x^4$.

\begin{figure}[t]
\centering
\begin{tikzpicture}[font=\small,
  plain/.style={minimum width=6.5mm,minimum height=4.2mm,inner sep=0pt},
  diag/.style={plain,fill=black!18},
  sub/.style={plain,draw=black!55,densely dotted},
  sup/.style={plain,fill=black!45,text=white},
  one/.style={plain,fill=black!18},
  two/.style={plain,fill=black!45,text=white},
  hdr/.style={inner sep=1pt,text=black!70}]
\node[hdr] at (0.00,0) {$k{=}0$}; \node[hdr] at (0.78,0) {1}; \node[hdr] at (1.56,0) {2}; \node[hdr] at (2.34,0) {3}; \node[hdr] at (3.12,0) {4}; \node[hdr] at (3.90,0) {5}; \node[hdr] at (4.68,0) {6}; \node[hdr] at (5.46,0) {7}; \node[hdr] at (6.24,0) {8}; \node[hdr] at (7.02,0) {9};
\node[hdr,anchor=east] at (-0.45,-0.52) {$n{=}1$}; \node[hdr,anchor=east] at (-0.45,-1.04) {2}; \node[hdr,anchor=east] at (-0.45,-1.56) {3}; \node[hdr,anchor=east] at (-0.45,-2.08) {4}; \node[hdr,anchor=east] at (-0.45,-2.60) {5}; \node[hdr,anchor=east] at (-0.45,-3.12) {6}; \node[hdr,anchor=east] at (-0.45,-3.64) {7};
\node[sub] at (0.00,-0.52) {1};
\node[diag] at (0.78,-0.52) {0};
\node[sup] at (1.56,-0.52) {0};
\node[plain] at (2.34,-0.52) {0};
\node[plain] at (3.12,-0.52) {0};
\node[plain] at (3.90,-0.52) {0};
\node[plain] at (4.68,-0.52) {0};
\node[plain] at (5.46,-0.52) {0};
\node[plain] at (6.24,-0.52) {0};
\node[plain] at (7.02,-0.52) {0};
\node[plain] at (0.00,-1.04) {1};
\node[sub] at (0.78,-1.04) {1};
\node[diag] at (1.56,-1.04) {0};
\node[sup] at (2.34,-1.04) {0};
\node[plain] at (3.12,-1.04) {0};
\node[plain] at (3.90,-1.04) {0};
\node[plain] at (4.68,-1.04) {0};
\node[plain] at (5.46,-1.04) {0};
\node[plain] at (6.24,-1.04) {0};
\node[plain] at (7.02,-1.04) {0};
\node[plain] at (0.00,-1.56) {1};
\node[plain] at (0.78,-1.56) {2};
\node[sub] at (1.56,-1.56) {2};
\node[diag] at (2.34,-1.56) {1};
\node[sup] at (3.12,-1.56) {0};
\node[plain] at (3.90,-1.56) {0};
\node[plain] at (4.68,-1.56) {0};
\node[plain] at (5.46,-1.56) {0};
\node[plain] at (6.24,-1.56) {0};
\node[plain] at (7.02,-1.56) {0};
\node[plain] at (0.00,-2.08) {1};
\node[plain] at (0.78,-2.08) {3};
\node[plain] at (1.56,-2.08) {5};
\node[sub] at (2.34,-2.08) {6};
\node[diag] at (3.12,-2.08) {5};
\node[sup] at (3.90,-2.08) {3};
\node[plain] at (4.68,-2.08) {1};
\node[plain] at (5.46,-2.08) {0};
\node[plain] at (6.24,-2.08) {0};
\node[plain] at (7.02,-2.08) {0};
\node[plain] at (0.00,-2.60) {1};
\node[plain] at (0.78,-2.60) {4};
\node[plain] at (1.56,-2.60) {9};
\node[plain] at (2.34,-2.60) {15};
\node[sub] at (3.12,-2.60) {20};
\node[diag] at (3.90,-2.60) {22};
\node[sup] at (4.68,-2.60) {20};
\node[plain] at (5.46,-2.60) {15};
\node[plain] at (6.24,-2.60) {9};
\node[plain] at (7.02,-2.60) {4};
\node[plain] at (0.00,-3.12) {1};
\node[plain] at (0.78,-3.12) {5};
\node[plain] at (1.56,-3.12) {14};
\node[plain] at (2.34,-3.12) {29};
\node[plain] at (3.12,-3.12) {49};
\node[sub] at (3.90,-3.12) {71};
\node[diag] at (4.68,-3.12) {90};
\node[sup] at (5.46,-3.12) {101};
\node[plain] at (6.24,-3.12) {101};
\node[plain] at (7.02,-3.12) {90};
\node[plain] at (0.00,-3.64) {1};
\node[plain] at (0.78,-3.64) {6};
\node[plain] at (1.56,-3.64) {20};
\node[plain] at (2.34,-3.64) {49};
\node[plain] at (3.12,-3.64) {98};
\node[plain] at (3.90,-3.64) {169};
\node[sub] at (4.68,-3.64) {259};
\node[diag] at (5.46,-3.64) {359};
\node[sup] at (6.24,-3.64) {455};
\node[plain] at (7.02,-3.64) {531};
\node[hdr,anchor=east] at (-0.45,-4.45) {$m$};
\node[hdr,anchor=east] at (-0.45,-5.0) {$p_{>3}(m)$};
\node[plain] at (0.00,-5.0) {1}; \node[hdr] at (0.00,-4.45) {0}; \node[plain] at (0.78,-5.0) {0}; \node[hdr] at (0.78,-4.45) {1}; \node[plain] at (1.56,-5.0) {0}; \node[hdr] at (1.56,-4.45) {2}; \node[plain] at (2.34,-5.0) {0}; \node[hdr] at (2.34,-4.45) {3}; \node[one] at (3.12,-5.0) {1}; \node[hdr] at (3.12,-4.45) {4}; \node[one] at (3.90,-5.0) {1}; \node[hdr] at (3.90,-4.45) {5}; \node[one] at (4.68,-5.0) {1}; \node[hdr] at (4.68,-4.45) {6}; \node[one] at (5.46,-5.0) {1}; \node[hdr] at (5.46,-4.45) {7}; \node[two] at (6.24,-5.0) {2}; \node[hdr] at (6.24,-4.45) {8};
\end{tikzpicture}
\caption{Top: $I_n(k)$ for $n\le7$, $k\le9$. Light cells are the main diagonal $k=n$ (giving $S_0$), dotted cells the
subdiagonal $k=n-1$ ($S_1$), dark cells the first superdiagonal $k=n+1$ ($S_{-1}$). Bottom: the coefficients
$p_{>n}(m)$ of $\prod_{j>n}(1-q^j)^{-1}$ for $n=3$. They vanish for $1\le m\le n$, equal $1$ for $n<m\le2n+1$
(one part), and first reach $2$ at $m=2n+2$, where $(n+1)+(n+1)$ appears. That single extra partition produces the
leading term of $P_i$.}
\label{fig:mech}
\end{figure}


\emph{Related work.} For $k\le n$ the numbers $I_n(k)$ are given by the Knuth--Netto formula (see \cite[Sec.~1]{CFS}
and the references there). Mennen \cite{M} extended a coefficient formula of the same kind to $k<2n+2$ by truncating
Euler's product; in the range $n\ge i-1$ this is the coefficient statement of Lemma~\ref{split} below. What is new
here is the generating-function identity (1) for every $i$, the formula for $P_i$, and its exact degree, which comes
from the one place where that truncation is not single-part.

\begin{theorem}\label{thm}
Conjecture~15 of \cite{CFS} is true. Write $F(q)=\prod_{j\ge1}(1-q^j)=\sum_s f_sq^s$, and for integers $n$ and $k$ let
$A(n,k)=[q^k]\,F(q)(1-q)^{-n}$. Then \textup{(1)} holds with
\[
P_i(x)=(1-x)^i\Bigl(\sum_{n=0}^{i-2}c_n x^{n+i}-\sum_{r=1}^{i}A(-r,i-r)\,x^{i-r}\Bigr)
+x^i\sum_{t=0}^{i-1}\Bigl(f_t(1-x)^{i-1}+\sum_{s=0}^{t-1}f_s\,x(1-x)^{i-t+s-1}\Bigr),
\]
where $c_n=I_n(n+i)-A(n,n+i)-\sum_{t=0}^{i-1}A(n,t)$. For $i\ge2$ the leading term of $P_i$ is $(-1)^ix^{3i-2}$.
\end{theorem}
In the formula for $P_i$, the first bracket holds the finitely many corrections from small $n$ and the polar part left
by Lagrange inversion; the last sum holds the contributions of single parts of $\prod_{j>n}(1-q^j)^{-1}$.

\section{Proof}
\emph{Strategy.} The Mahonian product $\prod_{j\le n}(1-q^j)/(1-q)$ is Euler's product $F(q)$ times $(1-q)^{-n}$,
corrected by the factor $\prod_{j>n}(1-q^j)^{-1}$. Without that factor, the diagonal generating functions would follow
from Lagrange inversion alone: for the subdiagonals this is the formula $S_i=(xC)^iS_0$, and for the superdiagonals
the same inversion, applied to a Laurent series, gives the main term $S_0C^{-i}$ plus a polynomial (Lemma~\ref{lagrange}).
The correcting factor contributes only through partitions into parts larger than $n$. Below degree $2n+2$ these are
single parts, whose contributions are polynomials in $n$ and so sum to rational functions with denominator
$(1-x)^{i}$ (Lemmas~\ref{split} and \ref{columns}). The first two-part partition, $(n+1)+(n+1)$, gives the top
coefficient of $P_i$.

Write $\phi(q)=1/(1-q)$ and $G_n(q)=\prod_{j>n}(1-q^j)^{-1}$, so that
\[
\sum_k I_n(k)q^k=\phi(q)^n F(q)\,G_n(q). \tag{2}
\]
Read (2) as: the Mahonian row is Euler's product spread out by $(1-q)^{-n}$, then corrected by $G_n$.
The coefficient of $q^m$ in $G_n$ is the number $p_{>n}(m)$ of partitions of $m$ into parts larger than $n$. Thus
$p_{>n}(0)=1$, $p_{>n}(m)=0$ for $1\le m\le n$, $p_{>n}(m)=1$ for $n<m\le 2n+1$ (the single part $m$), and
$p_{>n}(2n+2)=2$ (the partitions $2n+2$ and $(n+1)+(n+1)$).

\begin{lemma}\label{split}
Let $i\ge1$. For every $n\ge0$,
$I_n(n+i)=A(n,n+i)+\sum_{t=0}^{i-1}A(n,t)+c_n$, where $c_n=0$ for $n\ge i-1$ and $c_{i-2}=1$ when $i\ge2$.
\end{lemma}
\begin{proof}
By (2), $I_n(n+i)=\sum_{m=0}^{n+i}p_{>n}(m)\,A(n,n+i-m)$. The term $m=0$ is $A(n,n+i)$. Writing $t=n+i-m$, the terms
with $n<m\le n+i$ are $\sum_{t=0}^{i-1}p_{>n}(n+i-t)A(n,t)$, and the terms with $1\le m\le n$ vanish. If $n\ge i-1$ then
$n+i\le 2n+1$, so every $p_{>n}(n+i-t)$ equals $1$ and $c_n=0$. If $n=i-2$ then $n+i=2n+2$, the only coefficient
different from $1$ is $p_{>n}(2n+2)=2$ at $t=0$, and $c_{i-2}=A(i-2,0)=f_0=1$.
\end{proof}

The case $i=0$ of the same computation gives $I_n(n)=A(n,n)$ for all $n$, since $p_{>n}(m)=0$ for $1\le m\le n$.

\begin{lemma}\label{lagrange}
Let $w=xC(x)$, the power series with $w=x\phi(w)$. For every integer $i\ge0$,
\[
\sum_{n\ge0}A(n,n+i)\,x^n=x^{-i}C(x)^{-i}S_0(x)-\sum_{r=1}^{i}A(-r,i-r)\,x^{-r},
\]
where for $r\ge1$, $A(-r,k)=[q^k]F(q)(1-q)^r$.
\end{lemma}
So the main term of each superdiagonal is $S_0$ divided by a power of $xC(x)$, and what is lost in that division is a
polynomial in $1/x$ of degree at most $i$.
\begin{proof}
Lagrange inversion for Laurent series \cite[Theorem~2.1.1, (2.1.4)]{G}, with $R=\phi$ and $\psi(t)=t^{-i}F(t)$, gives for
every integer $n$
\[
[x^n]\,\frac{w^{-i}F(w)}{1-x\phi'(w)}=[t^n]\,t^{-i}F(t)\phi(t)^n=A(n,n+i).
\]
For $n<0$ the right side is $[q^{n+i}]F(q)(1-q)^{-n}$, which is zero unless $n\ge-i$. Hence
$w^{-i}F(w)/(1-x\phi'(w))=\sum_{n\ge0}A(n,n+i)x^n+\sum_{r=1}^{i}A(-r,i-r)x^{-r}$. With $i=0$ and $I_n(n)=A(n,n)$ this
reads $S_0=F(w)/(1-x\phi'(w))$, and since $w^{-i}=x^{-i}C^{-i}$ the general case follows.
\end{proof}

\begin{lemma}\label{columns}
For every $t\ge0$,
$\displaystyle\sum_{n\ge0}A(n,t)x^n=\frac{f_t}{1-x}+\sum_{s=0}^{t-1}\frac{f_s\,x}{(1-x)^{t-s+1}}$.
\end{lemma}
\begin{proof}
$\sum_{n\ge0}x^n\phi(q)^n=(1-q)/(1-x-q)$, whose coefficient of $q^u$ is $1/(1-x)$ for $u=0$ and
$(1-x)^{-u-1}-(1-x)^{-u}=x(1-x)^{-u-1}$ for $u\ge1$. Multiply by $F(q)$ and take the coefficient of $q^t$.
\end{proof}

\begin{proof}[Proof of Theorem~\ref{thm}]
Multiply Lemma~\ref{split} by $x^{n+i}$ and sum over $n\ge0$. Lemma~\ref{lagrange} turns the first term into
$C^{-i}S_0-\sum_{r=1}^{i}A(-r,i-r)x^{i-r}$, and Lemma~\ref{columns} turns the second into $x^i$ times a rational function
whose denominators divide $(1-x)^i$, because $t\le i-1$. Collecting terms gives (1) with the stated $P_i$, which has
integer coefficients because every $f_s$, $A(-r,k)$ and $c_n$ is an integer.

For the degree, the first sum in $P_i$ contributes degree at most $(i-1)+i=2i-1$ through the $A(-r,i-r)$ terms; each
summand of the last sum has degree at most $i-1$, so that part has degree at most $2i-1$; and
$(1-x)^i\sum_{n\le i-2}c_nx^{n+i}$ has degree at most $3i-2$, with coefficient $(-1)^ic_{i-2}=(-1)^i$ at $x^{3i-2}$. For
$i\ge2$ we have $3i-2>2i-1$, so $P_i$ has degree exactly $3i-2$. For $i=1$ the formula gives $P_1=-1+2x$.
\end{proof}

\begin{remark}
Only $c_{i-2}$ matters for the degree. In general
$c_n=\sum_{t=0}^{i-1}\bigl(p_{>n}(n+i-t)-1\bigr)A(n,t)$, a signed weighted sum (for instance $c_0=-1$ when $i=7$); the
terms with $n<i-2$ affect $P_i$ only below degree $3i-2$.
\end{remark}

\section*{Verification}
Independent exact checks (integer power series, no shared code with the proof's formula), over the stated finite ranges: the decomposition of
Lemmas~\ref{split} and \ref{lagrange} for $1\le i\le6$, and the explicit formula for $P_i$ against
$(1-x)^i\bigl(x^iS_{-i}-C^{-i}S_0\bigr)$ computed from Mahonian rows, for $1\le i\le 8$ to order $x^{40}$ or more. The
printed $P_1,\dots,P_4$ of \cite{CFS} are reproduced coefficient by coefficient (asserted in the check). A mutation suite plants ten false variants of the claims
(degree, sign, the value of $c_{i-2}$, exponents, a printed coefficient) and every one is rejected.

Lemmas~\ref{split} and \ref{columns}, the assembly of (1), the formula for $P_i$, its degree $3i-2$ and its leading
coefficient are formalised in Lean~4 with Mathlib. Lemma~\ref{lagrange} enters the formalisation as a stated
hypothesis citing \cite{G}; everything else is proved, using only the standard axioms. Code, scripts and the Lean
project: \url{https://github.com/jvvk/mathematics/tree/main/mahonian-superdiagonals}, archived at
doi:\href{https://doi.org/10.5281/zenodo.24536366}{10.5281/zenodo.24536366}.

\section*{Acknowledgement}
The proof was found with Claude (Anthropic, Claude Opus 5.5), which also wrote the verification scripts and, with a
subagent, the Lean formalisation. Codex (OpenAI) ran an independent literature search, an independent line-by-line
check of the proof and the second verification script. The author is responsible for the content.

\begin{thebibliography}{9}
\bibitem{CFS} A.~Claesson, A.~F.~Frankl\'in and E.~Steingr\'imsson, Permutations with few inversions,
\emph{Electron. J. Combin.} 30(4) (2023), \#P4.7, doi:10.37236/12075.
\bibitem{F} A.~F.~Frankl\'in, \emph{Enumerating score sequences and permutations by inversions and forbidden patterns},
PhD thesis, University of Iceland, 2026.
\bibitem{M} A.~Mennen, Mahonian triangle, unpublished note, 31 December 2010,
\url{https://www.alex.mennen.org/mahoniantri.pdf} (accessed 25 November 2026).
\bibitem{G} I.~M.~Gessel, Lagrange inversion, \emph{J. Combin. Theory Ser. A} 144 (2016), 212--249, arXiv:1609.05988.
\end{thebibliography}
\end{document}
